\scanpage{39}
le

\uline{Corollaire} Supposons que les $Z_i$ soient \uline{rares} dans $S$. Alors il existe un $x\in S$, avec $\dim\overline{\{x\}}=1$, tel que $x\notin\bigcup Z_i$.

Remplaçant $S$ par $T=\operatorname{Spec}V[[T_1,\ldots,T_r]]$, \deleted{\unsure{on voit}} on voit que les $T_i=\operatorname{Im}(Z_i\mathbin{-}T)$ sont rares, et il suffit de trouver un $x\in T$, avec $\dim\overline{\{x\}}=1$, tel que $x\notin\bigcup T_i$. (En effet, il existe alors un $y\in S$, au dessus de $x$, et on a nécessairement $\dim\overline{\{y\}}=1$).

On va montrer qu'il existe même un $x\in T$ provenant d'un $V$-hom
\[
V[[T_1,\ldots,T_r]]\mathbin{-}V.
\]
\deleted{(\unsure{continu})}

En effet, l'ens. de ces hom s'identifie à $\mathfrak m^r$ ($\mathfrak m=\operatorname{rad}(V)$). $\unsure{\mathcal H|\mathcal K}$ a une structure d'espace métrique complet. Montrons que pour tte partie fermée rare $Z$ de $T$, l'ens des $\varphi$ tels que $V(\operatorname{Ker}\varphi)\subset Z$ est rare. En effet, on peut supposer que \deleted{$Z$ est défini} $Z=V(f)$, \inserted{$f\in V[[T_1,\ldots,T_r]]$, $f\neq0$,} et $|V(\operatorname{Ker}\varphi)|\subset|V(f)|$ équivaut (car $V(\operatorname{Ker}\varphi)$ est intègre donc réduit) à
\[
V(\operatorname{Ker}\varphi)\subset V(f)\quad\text{i.e.}\quad f\in\operatorname{Ker}\varphi\quad\unsure{\text{i.e.}|\therefore}\quad\varphi(f)=0;
\]
i.e.
\[
f(\varphi_1,\ldots,\varphi_r)=0.
\]
Il est immédiat que c'est une partie fermée rare.

\mspage{38}
